Your install has roughly 48 instruments and 85 effects. This is a problem, not a feature.
Plugin browsing feels like work and is not work.

Here is the short list. Everything below is confirmed present on this machine.

\subsection{The six instruments to learn}\label{fl05_plugins-the-six-instruments-to-learn}

\textbf{FLEX.} Preset-driven, covers a huge range, and gets you a usable sound in ten seconds.
This is the one to reach for when you have an idea and do not want to lose it to sound
design. Free additional preset packs download from inside the plugin. Start here for pads,
plucks, keys and basses.

\textbf{FL Keys.} An electric piano. Play seventh and ninth chords on it and you are immediately
in deep house territory. It is not a great Rhodes emulation and that does not matter,
because a slightly cheap Rhodes is exactly what those records used.

\textbf{Transistor Bass.} FL\textquotesingle s Roland TB-303 emulation. This is the acid house instrument, and
having it bundled is genuinely lucky. Program a sixteenth-note pattern, turn resonance high,
sweep the cutoff, and use the slide and accent controls. Without slide and accent you have a
filtered bass; with them you have acid. See \texttt{THEORY/02\_chicago\_house.md} for why this matters
in Chicago specifically.

\textbf{3x Osc.} Three oscillators, a filter, two envelopes. Deliberately primitive, and the best
possible teaching instrument because there is nowhere to hide. Do the one-hour exercise at
the end of \texttt{THEORY/05\_sound\_design.md} in this plugin.

\textbf{Sytrus.} FM synthesis. Reach for it when you need bells, metallic tones, electric pianos
or aggressive basses that subtractive synths cannot produce. Harder to control. Learn it
third or fourth, not first.

\textbf{FPC.} A drum pad sampler. Useful once you are loading your own kits rather than dragging
individual samples into the Channel Rack.

Honourable mentions you can ignore for six months: \textbf{Harmor} (extremely powerful additive
and resynthesis engine, steep learning curve), \textbf{Slicex} (chopping breaks and vocals,
essential if you go toward footwork or jungle), \textbf{Drumaxx} (physically modelled drums),
\textbf{Morphine}, \textbf{Ogun}, \textbf{Sawer}, \textbf{PoiZone}, \textbf{Kepler}.

\subsection{The effects you will use every session}\label{fl05_plugins-the-effects-you-will-use-every-session}

\textbf{Fruity Parametric EQ 2.} The EQ. Highpass filters, surgical cuts, everything. You will
put this on more channels than any other plugin.

\textbf{Fruity Limiter.} Does three jobs: limiting on the master, compression (switch to the COMP
tab), and sidechain compression. The sidechain function is why it matters. See
\texttt{FL\_STUDIO/03\_first\_track.md}.

\textbf{Fruity Reeverb 2.} The reverb. Put it on a send, not on individual channels.

\textbf{Fruity Delay 3.} Tempo-synced delays, including the dotted-eighth setting that is standard
across dance music.

\textbf{Fruity Balance.} Just volume and pan, with no processing. Useful for trimming levels
before other plugins.

\textbf{Fruity Compressor.} A straightforward compressor for when you do not need the Limiter\textquotesingle s
extra features.

\subsection{The effects worth knowing about}\label{fl05_plugins-the-effects-worth-knowing-about}

\textbf{Maximus.} Multiband compression and limiting. This is a mastering-grade tool and it is
bundled. Overkill for now, but it is what you graduate to.

\textbf{Soundgoodizer.} A simplified front end for Maximus with one knob. It is a cheat, it is
slightly embarrassing, and it makes things sound better. Use it on the master while you are
learning, and stop using it once you understand what it is doing.

\textbf{Vintage Chorus.} Emulates the chorus circuit found on old Roland polysynths. Put it on a
pad and you are much closer to the Juno-60 sound that made "Can You Feel It" sound the way
it does.

\textbf{Gross Beat.} Time and volume manipulation with sixteen storable patterns. Two uses:
instant stutter and glitch effects, and as an alternative sidechain tool where you draw the
volume shape by hand. Widely used and instantly recognisable.

\textbf{Fruity Peak Controller.} Lets any signal control any parameter. The classic use is
sidechaining, but it can also make a filter open in time with the kick.

\textbf{Fruity Love Philter.} A multi-stage filter with envelope control. Good for the big filter
sweeps in breakdowns.

\textbf{Transient Processor.} Adjusts attack and sustain separately. Makes a soft kick punchy or
a sharp clap rounder, and it is more precise than compression for that job.

\textbf{Fruity Blood Overdrive} and \textbf{Fruity Fast Dist.} Distortion and saturation. A small
amount on a kick or bass adds harmonics that make the low end audible on small speakers.

\textbf{Fruity Stereo Enhancer.} Width control, and importantly the ability to force low
frequencies to mono, which you should be doing below about 150 Hz.

\textbf{Edison.} A full audio editor inside FL. Record, trim, reverse, clean up and export
samples. Reversing a crash cymbal for a transition is a two-minute job in here.

\textbf{Newtone.} Pitch and time correction for vocals and monophonic audio. If you ever record
singing, this is your tool.

\textbf{Vocodex.} A vocoder, for robot voices. Not core, but fun and historically appropriate.

\textbf{Wave Candy.} Spectrum analyser and oscilloscope. Put one on the master so you can see
what your low end is doing rather than guessing on speakers that cannot reproduce it.

\subsection{What you do not need to buy}\label{fl05_plugins-what-you-do-not-need-to-buy}

Nothing, for at least a year.

The plugins people will tell you to buy are Serum or Vital for wavetable synthesis, Kickstart
for sidechaining, and a good reverb. Vital is free and genuinely excellent, so if you want a
modern wavetable synth, get that. The other two duplicate things you already have.

Sample packs are a more defensible purchase than plugins, because good drum samples save
real time. But you have the entire 909, 808 and 707 palette bundled already, and those are
the correct sounds for what you are making.

\subsection{A rule for plugin browsing}\label{fl05_plugins-a-rule-for-plugin-browsing}

Pick one instrument per session and use only that one. Constraint produces finished tracks;
choice produces afternoons spent auditioning presets.
